\section{Implementation}\label{sec:impl}

The project is a Python package (\texttt{src/tan}, \LinesOfCode{} lines) built on NumPy and
SciPy, with \NumTests{} unit tests and continuous integration. \Cref{fig:arch} shows the
structure. All particle operations are vectorised over the particle array: no Python loop runs
over particles anywhere.

\begin{figure}[t]
  \centering
  \begin{tikzpicture}[
      box/.style={draw, rounded corners=2pt, align=center, font=\footnotesize,
                  minimum height=9mm, inner sep=3pt},
      arr/.style={-{Stealth[length=2mm]}, semithick}]
    \node[box] (data) {\texttt{data.py}\\ OpenTopography, reproject,\\ crop, manifest};
    \node[box, right=8mm of data] (dem) {\texttt{dem.py}\\ bilinear $h(\pv)$, $\nabla h$};
    \node[box, right=8mm of dem] (sim) {\texttt{simulate.py}\\ truth, INS, altimeters,\\ named random streams};
    \node[box, below=9mm of dem] (filters) {\texttt{filters/}\\ PF, RBPF, EKF,\\ dead reckoning, resampling};
    \node[box, right=8mm of filters] (runner) {\texttt{runner.py}\\ parallel, resumable\\ Monte Carlo};
    \node[box, right=8mm of runner] (metrics) {\texttt{metrics.py}, \texttt{pcrb.py}};
    \node[box, below=9mm of runner] (analysis) {\texttt{analysis/}\\ figures, tables,\\ \LaTeX{} macros};
    \draw[arr] (data) -- (dem);
    \draw[arr] (dem) -- (sim);
    \draw[arr] (dem) -- (filters);
    \draw[arr] (sim) -- (runner);
    \draw[arr] (filters) -- (runner);
    \draw[arr] (runner) -- (analysis);
    \draw[arr] (metrics) |- (analysis);
  \end{tikzpicture}
  \caption{Package structure and data flow. Every estimator implements the same interface
    (\texttt{initialize}, \texttt{predict}, \texttt{update}, \texttt{estimate}), so the runner
    treats them identically. The report reads only the outputs of \texttt{analysis/}.}
  \label{fig:arch}
\end{figure}

\subsection{The particle filter}
\Cref{lst:predict,lst:update} are the predict and update steps of \cref{eq:sir}, copied from the
source by a script so that they cannot drift from the code. The particles are an $N\times4$
array; the prediction samples all $N$ motion-model draws at once, each particle integrating the
INS velocity corrected by its own bias. The update evaluates the DEM at all particle positions
with one vectorised bilinear lookup (\texttt{scipy.ndimage.map\_coordinates}), adds the
log-likelihood to the log-weights and renormalises with log-sum-exp. Particles that leave the DEM
get zero weight rather than reading a clamped edge height, which would compare the measurement with
fake flat terrain. The same pass computes the predicted measurement spread for the NIS of
\cref{sec:monitor} and the log evidence.

\lstinputlisting[caption={Particle filter prediction (\texttt{ParticleFilter.predict}).},
  label={lst:predict}]{../generated/code_pf_predict.py}

\lstinputlisting[caption={Particle filter measurement update (\texttt{ParticleFilter.update}).},
  label={lst:update}]{../generated/code_pf_update.py}

Systematic resampling (\cref{lst:systematic}) needs one uniform number: evenly spaced pointers
are located in the cumulative weights by a binary search. Pinning the last cumulative weight to
exactly 1 and searching with \texttt{side="right"} guarantee that a particle with zero weight is
never selected, even when a pointer lands on a boundary; a unit test checks this for all four
schemes.

\lstinputlisting[caption={Systematic resampling.},label={lst:systematic}]{../generated/code_systematic.py}

The Rao-Blackwellised filter stores its particles in the same $N\times4$ array, with the Kalman
bias means in the last two columns, so the parent class's weighting and resampling are reused
unchanged; only the prediction (\cref{lst:rbpf}) and the covariance differ.

\lstinputlisting[caption={RBPF prediction with the shared Riccati recursion
  (\texttt{RaoBlackwellizedPF.predict}).},label={lst:rbpf}]{../generated/code_rbpf_predict.py}

\subsection{Simulation}\label{sec:sim}
The simulator is the ``world'': it knows the true trajectory and is the only component that
reads the DEM at the true position. A flight starts at a uniformly random point with a random
heading, flies at 40\,m/s and weaves $\pm25^\circ$ with a 120\,s period. The whole flight stays
$3\times2000$\,m inside the map --- three times the largest initial uncertainty used --- so that
neither the truth nor the prior reaches the edge; flights that would are redrawn. The altitude is
constant: the highest terrain within 1\,km of the track plus 300\,m clearance. INS velocities
and altimeter readings follow \cref{sec:models}; optional radar outliers, a spatially correlated
map-error field and a constant offset are added for Experiment~4.

Each source of randomness has its own named random stream per Monte Carlo run (flight, INS,
altimeters, outliers, initial guess, map error, and one per filter family), and every draw has a
fixed size and is scaled afterwards. Run $i$ is therefore identical in every configuration: a
larger $\sigma_0$ scales the \emph{same} initial-error direction, a higher outlier rate adds
outliers \emph{at the same instants plus new ones}, and so on. Unit tests check these invariances.

\subsection{Experiments and reproducibility}
Experiments are YAML files: a grid of parameters, a list of regions and estimator variants and a
number of runs. One task simulates a flight once and runs every estimator on it; tasks run in
parallel worker processes that share the DEMs in memory, and each writes a compressed file of
time series and per-run metrics. An interrupted sweep resumes where it stopped, and each results
folder records the git commit, the data-manifest checksum, the resolved configurations and the
installed package versions. Figures, tables and every number quoted in this report are produced
by scripts from these saved results (\texttt{make figures}); the data are traced by SHA-256
checksums of the downloaded and processed rasters.

Runtimes come from a separate benchmark --- one process pinned to a performance core, one maths
thread --- because inside the parallel sweeps the workers share memory bandwidth and each runs
several times slower. \Cref{tab:runtime} lists the cost per filter step on the hill-country DEM:
the particle filter with $N = 5\,000$ needs \RtPfFiveZeroZeroZero{}\,ms per step, three orders
of magnitude faster than the 1\,Hz altimeter rate. The RBPF costs the same as the plain filter
at equal $N$ (\RtRbpfFiveZeroZeroZero{}\,ms): its bias update is cheap and it draws half as many
random numbers.

\begin{table}[t]
  \centering\small
  \caption{Runtime per filter step (median of \RtReps{} repetitions of a 300-step flight on the
    hill-country DEM, single process, one performance core of the machine in
    \cref{app:repro}).}\label{tab:runtime}
  % Generated by scripts — do not edit.
\begin{tabular}{lrrr}
\toprule
Estimator & $N$ & Median (ms/step) & IQR (ms/step) \\
\midrule
Dead reckoning & -- & 0.009 & 0.009--0.010 \\
EKF & -- & 0.046 & 0.046--0.047 \\
PF & 500 & 0.215 & 0.202--0.217 \\
PF & 1\,000 & 0.283 & 0.274--0.291 \\
PF & 2\,000 & 0.436 & 0.416--0.456 \\
PF & 5\,000 & 0.832 & 0.825--0.874 \\
PF & 10\,000 & 1.511 & 1.506--1.578 \\
PF & 20\,000 & 2.877 & 2.871--2.891 \\
PF & 50\,000 & 7.022 & 7.013--7.118 \\
RBPF & 500 & 0.242 & 0.230--0.247 \\
RBPF & 1\,000 & 0.296 & 0.292--0.310 \\
RBPF & 2\,000 & 0.434 & 0.433--0.447 \\
RBPF & 5\,000 & 0.832 & 0.829--0.872 \\
RBPF & 10\,000 & 1.483 & 1.477--1.485 \\
RBPF & 20\,000 & 2.797 & 2.790--2.799 \\
PF, Student-t & 5\,000 & 0.849 & 0.847--0.858 \\
PF + monitor & 5\,000 & 0.851 & 0.847--0.856 \\
\bottomrule
\end{tabular}

\end{table}

\subsection{Validation}\label{sec:validation}
Because the outputs of a particle filter are random, ``it converges on a test case'' is weak
evidence of correctness. The decisive test uses a terrain for which the exact answer is known: on
a tilted plane $h(\pv) = 0.05\,p_x + 0.03\,p_y + 200$ the measurement is linear and bilinear
interpolation is exact, so the EKF \emph{is} the Kalman filter and computes the exact posterior.
A correct particle filter must agree with it up to Monte Carlo error. Over \ValSeeds{} seeds and
all 300 steps, the particle filter's estimate never deviates from the Kalman estimate by more than
\ValPFPosMax{} Kalman standard deviations in position and \ValPFBiasMax{} in bias
(\cref{tab:validation}); the RBPF, with four times fewer particles, stays within
\ValRBPFPosMax{} and \ValRBPFBiasMax{} --- a first glimpse of the variance reduction that
Rao-Blackwellisation promises. On the same plane the PCRB recursion reproduces the Kalman
covariance to a relative difference of \ValPcrbRelDiff{}. Further unit tests cover the resampling
schemes (unbiased offspring counts, zero-weight particles never selected, systematic offspring
equal to $\lfloor N w\rfloor$ or $\lceil N w \rceil$), the likelihoods (normalised densities,
outlier robustness), the true-noise model of the bound (against the simulator's empirical noise),
the divergence monitor (silent while tracking, recovers from a forced false lock), the metrics
(Wilson intervals, $\chi^2$ bounds, McNemar test against closed forms) and the runner (common
random numbers across configurations, resumption).

\begin{table}[t]
  \centering\small
  \caption{Particle filters against the exact Kalman filter on planar terrain ($\sigma_0 =
    100$\,m, \ValSeeds{} seeds, worst case over the 300-step flight).}\label{tab:validation}
  % Generated by scripts — do not edit.
\begin{tabular}{lrrrr}
\toprule
Filter & $N$ & max $|\hat x - \hat x_{KF}|/\sigma_{KF}$, position & same, bias & trace ratio $P/P_{KF}$ (position) \\
\midrule
PF & 20\,000 & 0.11--0.57 & 0.06--0.18 & 0.62--1.20 \\
RBPF & 5\,000 & 0.06--0.19 & 0.05--0.13 & 0.85--1.36 \\
\bottomrule
\end{tabular}

\end{table}

\Cref{tab:params} in the appendix lists all parameters.
